% !ZP title   = Newton's Three Laws of Motion
% !ZP type    = lecture
% !ZP topic   = Dynamics
% !ZP order   = 1
% !ZP summary = Inertia, F = ma, and action--reaction --- the foundation of classical mechanics.
% !ZP tags    = newton, inertia, momentum

\section*{Newton's Three Laws of Motion}

In \emph{kinematics} we described motion. In \textbf{dynamics} we explain
it. Almost all of classical mechanics rests on three laws.

\subsection*{First Law --- Inertia}
\begin{quote}
An object remains at rest, or moves in a straight line at constant
velocity, unless acted on by a net external force.
\end{quote}
\textbf{Inertia} is the tendency of matter to resist changes in motion;
mass is its measure.

\subsection*{Second Law --- $F = ma$}
\begin{equation}
  \vec{F}_{\text{net}} = m\vec{a}
\end{equation}
The same force produces a smaller acceleration on a larger mass.
A net force of $12\ \text{N}$ on a $3\ \text{kg}$ trolley gives
$a = 12/3 = 4\ \text{m/s}^{2}$.

\subsection*{Third Law --- Action and Reaction}
\begin{equation}
  \vec{F}_{A \to B} = -\,\vec{F}_{B \to A}
\end{equation}
These forces act on \textbf{different} objects, which is why they do not
cancel. A swimmer pushes water back; the water pushes the swimmer forward.

\subsection*{Summary}
\begin{center}
\begin{tabular}{lll}
\textbf{Law} & \textbf{Idea} & \textbf{Relation} \\
First  & Inertia            & $\vec{F}_{\text{net}}=0 \Rightarrow$ constant $\vec{v}$ \\
Second & Force--acceleration & $\vec{F}_{\text{net}} = m\vec{a}$ \\
Third  & Interaction pairs  & $\vec{F}_{A\to B} = -\vec{F}_{B\to A}$ \\
\end{tabular}
\end{center}
